\section{Conclusion}

We introduce obfuscation testing, a novel methodology for validating whether large language models detect structural patterns through causal reasoning rather than temporal memorization. By stripping all temporal and identity context from financial data while preserving mechanical relationships, we create a rigorous test of LLM reasoning capabilities independent of training data leakage.

Applied to dealer hedging constraints in options markets, LLMs achieve 71.5\% detection rate using unbiased prompts across 242 trading days in 2024 (95.6\% of trading days), significantly exceeding the 60\% mechanical threshold. More critically, a \textbf{raw chain validation} removing all pre-calculated metrics achieves 92.3\% detection—outperforming the GEX-assisted baseline by 30.8 percentage points on identical test dates. This result provides definitive evidence that the LLM functions as a \textbf{structural analyst}, reconstructing dealer gamma constraints from strike-level open interest distribution rather than pattern-matching on our ``\$-32B GEX'' calculations.

\subsection{Twin Pillars of Structural Reasoning}

Two findings provide converging evidence for genuine mechanism identification:

\textbf{Pillar 1: Raw Chain Superiority.} When we remove pre-calculated GEX entirely, the LLM's detection \textit{improves} rather than collapses. The 92.3\% raw chain performance (vs. 61.5\% GEX-assisted baseline) demonstrates that: (1) the LLM is not reading our metric labels, (2) it is performing the gamma integration itself from strike-by-strike OI patterns, and (3) it leverages richer contextual information (asymmetrical concentration, IV skew, distance from ATM) that scalar metrics necessarily discard. This positions LLMs as tools capable of extracting signal from data dimensionality that traditional parametric approaches reduce through lossy compression.

\textbf{Pillar 2: Inverse P-Hacking and Alpha Orthogonality.} Detection days exhibit 33\% \textit{lower} intraday range expansion than baseline (21.6\% vs. 32.0\%, p = 0.03), demonstrating the LLM detects \textit{volatility suppression} rather than chasing high-volatility outcomes. Unlike typical data mining exercises that search for high-volatility anomalies, our approach validates robustness by detecting the \textit{absence} of volatility when dealer constraints bind. Combined with 91.2\% materialization accuracy, this proves pattern-specific detection rather than correlation fishing.

Critically, the divergence between stable detection (68-74\% across quarters) and declining economic profitability (Sharpe 1.8 $\rightarrow$ 0.1) establishes these signals as \textbf{risk management indicators rather than alpha generators}. The lack of trading alpha is a feature, not a limitation: signals orthogonal to momentum and mean-reversion strategies provide diversification value for risk models precisely because they capture structural constraints rather than exploitable price patterns.

\subsection{Comprehensive Defense Framework}

We validate structural reasoning through five independent empirical tests, each addressing a specific methodological critique:

\begin{enumerate}
    \item \textbf{Sensitivity vs. Guessing}: Non-detection days exhibited 3.72$\times$ lower mean GEX magnitude (\$5.4B vs. \$20.1B) and significantly weaker signal concentration (p < 0.0001, Cohen's d = 0.68). The LLM requires concentrated structural constraints, refuting base rate guessing.

    \item \textbf{Inverse P-Hacking}: Detection days show 33\% lower range expansion than materialization days (p = 0.03). The LLM detects suppression, not volatility chasing.

    \item \textbf{Not Profit-Chasing}: Confidence increased (79.5 $\rightarrow$ 80.7) Q1-Q4 2024 while net alpha collapsed (Sharpe 1.8 $\rightarrow$ 0.1). The LLM identifies mechanisms independent of profitability.

    \item \textbf{EOD Validity}: Statistical baseline achieves AUC = 0.681 predicting T+1 outcomes from EOD features. GEX normalized by OI (p = 0.0075) is the only significant predictor, validating EOD structural snapshots.

    \item \textbf{Structural Analyst}: Raw chain (92.3\%) outperforms GEX-assisted (61.5\%) baseline by 30.8pp. The LLM reconstructs dealer constraints from first principles, not metric pattern-matching.
\end{enumerate}

Together, these five defenses establish that the LLM operates as a genuine structural analyst, identifying forced market actions from dispersed dealer positioning constraints.

\subsection{Contributions and Implications}

The primary contribution is the validation methodology itself. Obfuscation testing provides a systematic framework for distinguishing AI reasoning from memorization, applicable beyond finance to any domain with structural constraints. The WHO$\rightarrow$WHOM$\rightarrow$WHAT causal framework proves effective for forcing explicit mechanical attribution, preventing vague pattern claims and requiring specific identification of economic actors, affected parties, and forced actions.

The raw chain validation introduces a novel robustness test: \textit{removing} analytical scaffolding to test whether models genuinely perform structural reasoning. This approach inverts traditional ablation studies—rather than adding complexity, we strip pre-calculated features to verify the model reconstructs them independently. The 30.8pp performance improvement suggests that scalar metrics like net GEX may actually obscure structural signals visible in the full strike-level data. This has immediate implications for quantitative finance: LLMs may excel at pattern recognition from unstructured market data that parametric models necessarily reduce through dimensionality compression.

The inverse p-hacking finding establishes a rigorous robustness standard for financial ML research. Rather than seeking high-volatility outcomes (the typical data mining bias), the LLM detects \textit{suppression}—days when dealer hedging mechanically dampens realized volatility below the predicted range. This asymmetry provides strong evidence for mechanism-specific reasoning that cannot result from correlation fishing. For practitioners, this establishes the detected signal as a risk management input: structural constraints that suppress volatility are precisely the regime information risk models require for accurate variance forecasting.

As LLMs increasingly influence decision-making in regulated markets, rigorous validation of their capabilities becomes essential. Our framework—temporal obfuscation, causal attribution requirements, outcome verification, and raw chain robustness—establishes a reproducible standard for evaluating AI structural reasoning in financial applications.

\subsection{Future Directions}

Future work should validate detection across multiple asset classes and market regimes to establish generalizability. Testing on individual equities may reveal positive gamma regimes absent in our SPY sample, enabling traditional regime comparison analysis. Extending the temporal scope to crisis periods (2020 COVID crash, 2022 volatility spike) would test whether detection persists under extreme market stress.

Intraday analysis with high-frequency data could reveal microstructure patterns invisible at daily granularity, particularly for 0DTE options where dealer rehedging occurs within hours. The raw chain methodology could be extended to other structural patterns: order flow toxicity, quote stuffing, or liquidity provision constraints—any pattern where participants face forced actions under specific conditions.

Cross-model validation (GPT-4o-mini, GPT-o3-mini, Claude Sonnet) could distinguish model-specific artifacts from robust pattern detection through consensus mechanisms. The consistency of detection across models would further validate structural reasoning over memorization of training corpus artifacts.

We make our code and validation framework publicly available at https://github.com/iAmGiG/gex-llm-patterns to encourage replication and extension. The complete pipeline—including data obfuscation, LLM prompting templates, raw chain validation, outcome verification, and statistical testing—is provided with documentation enabling researchers to apply this methodology to other financial patterns or non-financial domains with structural constraints. As markets evolve and new constraints emerge, the ability to rigorously validate AI understanding of structural patterns will prove increasingly valuable for researchers, practitioners, and regulators seeking to deploy trustworthy AI systems in financial markets.
